The representative universal-property theorem is
\wikilink{s13--b006}{the abelianization theorem}. The quotient map goes
to $G/[G,G]$, as specified in the existing source correction.

\begin{proposition}\label{abelianization-intersection}
The commutator subgroup $[G,G]$ equals the intersection of all normal
subgroups $N$ for which $G/N$ is abelian.
\end{proposition}
\begin{proof}\label{abelianization-intersection-proof}
If $G/N$ is abelian, every commutator maps to the identity, so
$[G,G]\subseteq N$. Conversely $[G,G]$ is itself normal with abelian
quotient by \wikilink{s13--b004}{the commutator proposition}.
It is therefore one member of the family being intersected, proving
the opposite inclusion.
\end{proof}
\begin{remark}\label{abelianization-maps}
Consequently every map from $G$ to an abelian group kills this
intersection and factors uniquely through the quotient map.
The factorization, rather than a map into the commutator subgroup,
is what characterizes abelianization.
\end{remark}
